\section{\texorpdfstring{Computing Core Values over Static \cpi}{Computing Core Values over Static CPI}}\label{sec:hindex}

Both algorithms in this section are written over a single high-level operator: the number of $s$-cliques containing a vertex inside a given vertex set.
%
Section~\ref{sec:setsupport} names this operator $\setSup_R(v)$ and shows that the static \cpi{} evaluates it without enumerating cliques.
%
Section~\ref{sec:localh} casts \spin{} (Algorithm~\ref{alg:localH}) as a direct fixed-point iteration that queries $\setSup_R$ on threshold-induced sets.
%
Section~\ref{sec:peelh} derives \spinStar{} (Algorithm~\ref{alg:peelH}) by maintaining $\setSup_R$ on a residual set under single-vertex deletions; Section~\ref{sec:changeonly-impl} gives the static-\cpi{} realisation of that maintenance.

\subsection{Set support over the static \cpi}\label{sec:setsupport}

\begin{definition}[Set support]\label{def:set-support}
For a vertex set $R\subseteq V$ and a vertex $v\in R$, the \emph{set support} of $v$ in $R$ is
\[
  \setSup_R(v) \;=\; \bigl|\{S\subseteq R : v\in S,\ |S|=s,\ S\text{ is a clique in } G[R]\}\bigr|.
\]
The full-graph case recovers $s$-degree: $\setSup_V(v)=\sdeg(v)$.
\end{definition}

Both algorithms below are written over $\setSup_R$ alone.
%
\spin{} queries it on threshold-induced sets that change every pass; \spinStar{} maintains it on a residual set that shrinks by one vertex per pop.

The static \cpi{} evaluates $\setSup_R(v)$ as a sum of per-path contributions: every encoded $s$-clique sits on a unique path $\TPath$, and that clique lies inside $G[R]$ iff $\Vh(\TPath)\subseteq R$ and its $\eta_{\TPath}$ pivots are drawn from $\Vp(\TPath)\cap R$.

\begin{lemma}[\cpi{} evaluation of set support]\label{lem:set-support-eval}
For any $R\subseteq V$ and $v\in R$,
\[
  \setSup_R(v) \;=\; \sum_{\TPath\ni v} \mathbf{1}[\Vh(\TPath)\subseteq R]\cdot\pathContrib(\TPath,v,|\Vp(\TPath)\cap R|).
\]
\end{lemma}

\begin{proof}
Every $s$-clique containing $v$ is encoded by a unique path $\TPath$ as $\Vh(\TPath)\cup P'$ with $P'\subseteq\Vp(\TPath)$ and $|P'|=\eta_{\TPath}$ (Definition~\ref{def:path}).
%
Such a clique lies in $G[R]$ iff $\Vh(\TPath)\subseteq R$ and $P'\subseteq\Vp(\TPath)\cap R$; conditioned on the hold inclusion, the count of qualifying $P'$ is $\pathContrib(\TPath,v,|\Vp(\TPath)\cap R|)$ by Definition~\ref{def:path-contrib}.
%
Summing over paths through $v$ gives the stated total.
%
\qed
\end{proof}

Sorting each scanned path's vertex values once and reading the per-path contribution off prefix sums computes $\setSup_R(v)$ at per-call cost $O(w_v\log L_{\max})$, with $w_v=\sum_{\TPath\ni v}(|\Vh(\TPath)|+|\Vp(\TPath)|)$ and $L_{\max}=\max_{\TPath}(|\Vh(\TPath)|+|\Vp(\TPath)|)$.
%
The call reads the static \cpi{} and the set $R$, and it mutates no path.

\subsection{\texorpdfstring{\spin{}: direct fixed-point specification}{SPIN: direct fixed-point specification}}\label{sec:localh}

Given a value vector $h$, define the threshold-induced set
\[
  R_h(k) \;=\; \{u\in V : h[u]\ge k\}.
\]
The one-vertex update is the largest $k$ at which $v$ has enough $s$-cliques inside $R_h(k)$:
\[
  \Phi_h(v) \;=\; \max\bigl\{k\ge 0 \;:\; \setSup_{R_h(k)}(v)\ge k\bigr\}.
\]
\spin{} (Algorithm~\ref{alg:localH}) initialises $h[v]\gets\setSup_V(v)=\sdeg(v)$, then repeatedly applies $\Phi_h$ to every vertex until $h$ no longer changes.
%
The only mutable data are the current and next value vectors; the initial vector dominates the final core value pointwise, and each pass can only lower entries.
%
The converged vector is the canonical semantic anchor for any static-\cpi{} core-value solver.

\begin{algorithm}[t]
\algfontsize
\caption{\spin{}: direct fixed-point iteration over set support.}
\label{alg:localH}
\KwIn{Graph $G$, clique size $s$, \cpi{} $\mathcal{P}$.}
\KwOut{$\kappa_s(v)$ for every $v\in V$.}

\lForEach{$v\in V$}{$h[v]\gets\setSup_V(v)$}
$\mathrm{changed}\gets\textsf{true}$\;
\While{$\mathrm{changed}$}{
    $\mathrm{changed}\gets\textsf{false}$;\quad $h_{\mathrm{new}}\gets h$\;
    \ForEach{$v\in V$ \emph{in a fixed order}}{
        $h_{\mathrm{new}}[v]\gets\max\bigl\{k\le h[v] : \setSup_{R_h(k)}(v)\ge k\bigr\}$\tcp*{Lemma~\ref{lem:set-support-eval}}
        \lIf{$h_{\mathrm{new}}[v]\ne h[v]$}{$\mathrm{changed}\gets\textsf{true}$}
    }
    $h\gets h_{\mathrm{new}}$\;
}
\lForEach{$v\in V$}{$\kappa_s(v)\gets h[v]$}
\Return{$\kappa_s$}
\end{algorithm}

\begin{example}\label{ex:hop}
Consider the running example of Figure~\ref{fig:overview} (Table~\ref{tab:cpi-leaves}).
%
Initialise $h[v]\gets\sdeg(v)$ via Lemma~\ref{lem:supp}: for instance, $h[v_4]\!=\!6$, summing $2{+}1{+}2$ from $\TPath_3,\TPath_4,\TPath_5$ as pivot and $1$ from $\TPath_6$ as hold.
%
The first $\Phi_h$ evaluation at $v_4$ tests, for descending $k$, whether $\setSup_{R_h(k)}(v_4)\ge k$ using Lemma~\ref{lem:set-support-eval}; the largest passing $k$ is $3$, so $h[v_4]$ drops $6\!\to\!3$.
%
Table~\ref{tab:spinex} shows the full trajectory: iteration~$1$ drops the six $K_4$ vertices to $3$, and iteration~$2$ reproduces iteration~$1$, certifying convergence at $h^\star\!=\!\kappa_3$.
\end{example}

\begin{table}[t]
\centering
\caption{\spin{} trace on the running example of Figure~\ref{fig:overview}.}
\label{tab:spinex}
\tiny
\setlength{\tabcolsep}{4pt}
\begin{tabular}{l|cccccccccc}
\toprule
\multirow{2}{*}{iteration $t$} & \multicolumn{10}{c}{$h$-value $h^{(t)}[v]$ per vertex} \\
\cmidrule(lr){2-11}
                                 & $v_1$ & $v_2$ & $v_3$ & $v_4$ & $v_5$ & $v_6$ & $v_7$ & $v_8$ & $v_9$ & $v_{10}$ \\
\midrule
$0$ (initial $\sdeg$)            & $3$ & $4$ & $6$ & $6$ & $3$ & $4$ & $1$ & $1$ & $1$ & $1$ \\
$1$                              & $3$ & $3$ & $3$ & $3$ & $3$ & $3$ & $1$ & $1$ & $1$ & $1$ \\
$2$ (converged $=\kappa_3$)    & $\mathbf{3}$ & $\mathbf{3}$ & $\mathbf{3}$ & $\mathbf{3}$ & $\mathbf{3}$ & $\mathbf{3}$ & $\mathbf{1}$ & $\mathbf{1}$ & $\mathbf{1}$ & $\mathbf{1}$ \\
\bottomrule
\end{tabular}
\end{table}

\begin{lemma}[\spin{} Convergence]\label{lem:hindex-convergence}
Initialised at $h^{(0)}[v]=\sdeg(v)$, Algorithm~\ref{alg:localH} terminates in finitely many passes with $h[v]=\kappa_s(v)$ for every $v\in V$, and the sequence $h^{(t)}[v]$ is monotonically non-increasing.
\end{lemma}

\begin{proof}
By the operator definition, $h^{(t+1)}[v]\!\ge\!k$ iff $v$ has at least $k$ incident $s$-cliques inside $R_{h^{(t)}}(k)$, so $\sdeg(v)$ is the largest possible value and the sequence is monotone non-increasing on the non-negative integers.
%
The sequence is bounded below by $0$ and decreases by at least one in any non-converged pass, so the iteration terminates and a pass with no entry change certifies the converged vector $h^\star$.
%
At $h^\star$, let $S_k=R_{h^\star}(k)$, and let $C\subseteq S_k$ be $v$'s $s$-connected component within $S_k$.
%
Every $v\!\in\!S_k$ has $\setSup_{S_k}(v)\ge k$, and each such clique places $v$ in $C$ (the clique itself provides the link), giving Conditions~(i) and (ii) of Definition~\ref{def:nucleus}.
%
For (iii), suppose for contradiction that some $w\!\notin\!C$ could be added to $C$ while preserving (i) and (ii); then $w$ has at least $k$ $s$-cliques entirely inside $C\cup\{w\}\subseteq S_k\cup\{w\}$, so $\setSup_{S_k\cup\{w\}}(w)\ge k$ and hence $h^\star[w]\ge k$, placing $w\in S_k$.
%
Moreover each such clique links $w$ to $C$-vertices, so $w\in C$, a contradiction.
%
Hence $C$ is a $k$-$(1,s)$-nucleus, giving $\kappa_s(v)\!\ge\!k$.
%
Conversely, if $v$ lies in a $k$-$(1,s)$-nucleus $C'$, we show by induction on $t$ that $h^{(t)}[u]\!\ge\!k$ for every $u\!\in\!C'$.
%
The base case $t\!=\!0$ holds since each $u\!\in\!C'$ is in at least $k$ $s$-cliques inside $C'$, so $h^{(0)}[u]\!=\!\sdeg(u)\!\ge\!k$.
%
For the inductive step, $C'\subseteq R_{h^{(t)}}(k)$ implies $\setSup_{R_{h^{(t)}}(k)}(u)\!\ge\!\setSup_{C'}(u)\!\ge\!k$, so $h^{(t+1)}[u]\!\ge\!k$.
%
Therefore $h^\star[v]\!\ge\!k$, and hence $h^\star[v]=\kappa_s(v)$.
%
\qed
\end{proof}

\begin{theorem}[\spin{} Complexity]\label{thm:localh-complexity}
Algorithm~\ref{alg:localH} computes $\kappa_s$ in $O(T\cdot\Sig\cdot L_{\max}\cdot\log L_{\max})$ work, where $\Sig$ is the static \cpi{}'s total vertex--path incidence count (Section~\ref{sec:prelim}), $T$ is the number of passes until convergence, and $L_{\max}$ is the maximum path size.
\end{theorem}

\begin{proof}
Each pass evaluates $\Phi_h(v)$ for every $v\in V$ at per-call cost $O(w_v\log L_{\max})$ (Section~\ref{sec:setsupport}).
%
Summing, $\sum_v w_v=\sum_{\TPath}|\TPath|^2\le L_{\max}\cdot\Sig$, so each of the $T$ passes costs $O(\Sig\cdot L_{\max}\cdot\log L_{\max})$.
%
\qed
\end{proof}

\begin{remark}[Inefficiencies of \spin{}]\label{rem:spin-ineff}
\spin{} pays two kinds of redundant work.
%
First, it re-evaluates $\Phi_h$ on every vertex per pass for up to $T$ passes, so a vertex whose final value is reached at pass~$1$ is still refreshed in passes $2,3,\ldots,T$.
%
Second, each evaluation recomputes $\setSup_{R_h(k)}$ from scratch over all paths through $v$ at every threshold $k$, repeating the binomial counts the previous pass already produced.
\end{remark}

\subsection{\texorpdfstring{\spinStar{}: change-only peeling}{SPIN-star: change-only peeling}}\label{sec:peelh}

\spinStar{} replaces \spin{}'s repeated full evaluations with \emph{change-only peeling} on a single mutable set.
%
The set is the \emph{residual set} $R$, initialised to $V$ and shrunk by one vertex per pop in deletion order.
%
Instead of recomputing $\setSup_R$ from scratch, \spinStar{} maintains $\supArr[v]=\setSup_R(v)$ for every unfinalised $v$ and updates it under each deletion.

\stitle{Change-only delta.}
When $R$ changes to $R\setminus\{x\}$ by deleting $x$, only vertices $u$ that share an encoded $s$-clique with $x$ can lose support; for every such $u$, the support drop is
\[
  \Delta_x(u) \;=\; \setSup_R(u) - \setSup_{R\setminus\{x\}}(u).
\]
\spinStar{} updates only the affected $\supArr[u]\gets\supArr[u]-\Delta_x(u)$, rather than recomputing $\setSup$ from scratch on the new residual set.

\stitle{Peel-order schedule.}
The priority-queue key is $\keyArr[v]=\max(k_{\mathrm{cur}},\supArr[v])$, where $k_{\mathrm{cur}}$ is the current peel level.
%
The clip keeps the extracted-key sequence monotone non-decreasing, so each pop's $k$ equals the popped vertex's final $\kappa_s$.
%
The path-level realisation of $\Delta_x(\cdot)$ on the static \cpi{} is given in Section~\ref{sec:changeonly-impl}.

\begin{algorithm}[t]
\algfontsize
\caption{\spinStar{}: change-only peeling over set support.}
\label{alg:peelH}
\KwIn{Graph $G$, clique size $s$, \cpi{} $\mathcal{P}$.}
\KwOut{$\kappa_s(v)$ for every $v\in V$.}

$R\gets V$\;
\lForEach{$v\in V$}{$\supArr[v]\gets\setSup_R(v)$;\quad $\keyArr[v]\gets\supArr[v]$}
$Q\gets$ min-priority queue with key $\keyArr[\cdot]$; insert each $v\in V$\;
\While{$R\ne\emptyset$}{
    $(v,k)\gets Q.\textsc{extract-min}()$\;
    $\kappa_s(v)\gets k$;\quad $R\gets R\setminus\{v\}$\;
    \ForEach{$u\in R$ \emph{with} $\Delta_v(u)\ne 0$}{
        $\supArr[u]\gets\supArr[u]-\Delta_v(u)$\tcp*{Lemma~\ref{lem:changeonly-cpi}, $O(1)$ per path}
        $y\gets\max\{k,\,\supArr[u]\}$\;
        \lIf{$y<\keyArr[u]$}{$Q.\textsc{decrease-key}(u,y)$;\quad $\keyArr[u]\gets y$}
    }
}
\Return{$\kappa_s$}
\end{algorithm}

Algorithm~\ref{alg:peelH} initialises $R\!=\!V$, $\supArr[v]\!=\!\setSup_V(v)\!=\!\sdeg(v)$, and seeds the priority queue with $\keyArr[\cdot]$ (Lines~1-3).
%
The peel loop (Lines~4-9) extracts the minimum-key vertex $v$ at peel level $k$, finalises $\kappa_s(v)\!\gets\!k$, removes $v$ from $R$, and propagates the change-only delta $\Delta_v$ to every affected unfinalised $u$ (Line~7); each $\supArr[u]$ is decremented (Line~8) and $u$'s priority-queue key is lifted via $\textsc{decrease-key}$ when the clipped value drops (Line~9).
%
Because $\keyArr[u]\ge k$ for every unfinalised $u$ and $k$ is non-decreasing across extractions, the extracted-key sequence is monotone, so each pop assigns its final $\kappa_s$ value~\cite{coreAlg}.

\begin{example}\label{ex:peelh}
Same setup as Example~\ref{ex:hop}.
%
At iteration~$1$ the four pendant-triangle vertices $v_7,v_8,v_9,v_{10}$ tie at minimum key $1$; suppose $v_8$ pops first at $k\!=\!1$.
%
Removing $v_8$ from $R$ retires $\TPath_1$ (since $\Vp(\TPath_1)\cap R$ loses its only spare pivot, killing the path); the affected co-path vertices are $v_7$ (hold of $\TPath_1$, loses $\binom{2}{2}\!=\!1$) and $v_6$ (pivot of $\TPath_1$, loses $\binom{1}{1}\!=\!1$).
%
At iteration~$7$ the extracted key reaches $k\!=\!3$ and $v_4$ pops, sending change-only deltas through the two live paths $\TPath_5$ and $\TPath_6$ as in Lemma~\ref{lem:changeonly-cpi}.
%
After ten pops, $\kappa_3(v_i)\!=\!3$ for $i\!\in\!\{1,\ldots,6\}$ and $\kappa_3(v_i)\!=\!1$ for $i\!\in\!\{7,\ldots,10\}$; Table~\ref{tab:peel} traces $\supArr$ across the ten iterations.
\end{example}

\begin{lemma}[\spinStar{} Equivalence]\label{lem:equivalence}
Algorithm~\ref{alg:peelH} returns the same $\kappa_s$ values as Algorithm~\ref{alg:localH}.
\end{lemma}

\begin{proof}
We establish the \emph{residual-support invariant}: at every iteration of the peel loop, $\supArr[u]=\setSup_R(u)$ for every $u\in R$.

\smallskip
\noindent\emph{(i) Initialisation.}
At $t=0$, $R=V$ and $\supArr[v]\!\gets\!\setSup_V(v)$ directly (Definition~\ref{def:set-support}).

\smallskip
\noindent\emph{(ii) Inductive step.}
Assume the invariant holds when $v$ is popped.
%
Removing $v$ from $R$ changes $\setSup_R(u)$ to $\setSup_{R\setminus\{v\}}(u)=\setSup_R(u)-\Delta_v(u)$ for every other $u\in R$, by Definition~\ref{def:set-support}.
%
Algorithm~\ref{alg:peelH} subtracts exactly $\Delta_v(u)$ from $\supArr[u]$ (Line~8), restricted to $u$ with $\Delta_v(u)\!\ne\!0$ (Line~7) since other vertices' supports are unchanged.
%
After the iteration, $\supArr[u]=\setSup_{R\setminus\{v\}}(u)$ on the new residual, preserving the invariant.

\smallskip
\noindent\emph{(iii) Peel correctness.}
Algorithm~\ref{alg:peelH} is the standard deletion-order peeling template of $(1,s)$-core decomposition~\cite{coreAlg,nucleusSariyuce14}: at each iteration the unfinalised vertex of minimum $\max\{k_{\mathrm{cur}},\supArr[u]\}$ is finalised, and its set-support contribution is removed from every other vertex.
%
The clip $\keyArr[u]\gets\max\{k,\supArr[u]\}$ (Line~9) does not alter $\supArr[u]$; it enforces non-decreasing extraction by lifting the queue key when the residual support has already fallen below the current peel level.
%
By the standard correctness of counting-based core decomposition for $(r,s)$-nucleus~\cite{nucleusSariyuce14}, each popped vertex's $k$ equals its $(1,s)$-nucleus core number.

\smallskip
By Lemma~\ref{lem:hindex-convergence}, this $\kappa_s$ vector is the same fixed point reached by \spin{}.
%
\qed
\end{proof}

\begin{theorem}[\spinStar{} Complexity]\label{thm:peelh-complexity}
Algorithm~\ref{alg:peelH} runs in $O(\Sig+(U+n)\log n)$ time, where $\Sig$ is the total vertex--path incidence count of the static \cpi{} and $U$ is the total number of $\supArr$ updates over all pops.
\end{theorem}

\begin{proof}
Initialisation evaluates $\setSup_V$ for every $v$ at total cost $O(n+\Sig)$ (Lemma~\ref{lem:set-support-eval} amortised, since each vertex--path incidence is touched once), plus $n$ priority-queue insertions at $O(\log n)$ each.
%
Each subsequent pop walks the alive paths through the popped vertex once, contributing $O(\Sig_v)$ work for delta propagation (Lemma~\ref{lem:changeonly-cpi}); summing over the $n$ pops gives $O(\Sig)$ total.
%
The $n$ extractions plus $U$ decrease-keys contribute $O((U+n)\log n)$.
%
The total time is therefore $O(\Sig+(U+n)\log n)$.
%
\qed
\end{proof}

\begin{corollary}[Worst-case bound on $U$]\label{cor:U-bound}
$U\le L_{\max}\cdot\Sig$, where $L_{\max}=\max_{\TPath}(|\Vh(\TPath)|+|\Vp(\TPath)|)$.
%
Hence Algorithm~\ref{alg:peelH} runs in $O(\Sig\cdot L_{\max}\cdot\log n)$ time in the worst case.
\end{corollary}

\begin{proof}
For each path $\TPath$ and each $u\in\Vh(\TPath)\cup\Vp(\TPath)$, $u$ is updated at most once per pop of another vertex on $\TPath$, i.e.\ at most $|\TPath|-1$ times.
%
Summing over $(\TPath,u)$ pairs gives $U\le\sum_{\TPath}|\TPath|\cdot(|\TPath|-1)\le L_{\max}\cdot\Sig$.
%
Substituting into Theorem~\ref{thm:peelh-complexity} yields the worst-case bound.
%
\qed
\end{proof}

\begin{table}[t]
\centering
\caption{\spinStar{} trace on the running example of Figure~\ref{fig:overview}.}
\label{tab:peel}
\tiny
\setlength{\tabcolsep}{4pt}
\begin{tabular}{c|cccccccccc}
\toprule
\multirow{2}{*}{iteration $t$} & \multicolumn{10}{c}{residual support $\supArr^{(t)}[v]=\setSup_R(v)$ per vertex \quad (bold $=\kappa_3(v)$ at pop)} \\
\cmidrule(lr){2-11}
                                 & $v_1$ & $v_2$ & $v_3$ & $v_4$ & $v_5$ & $v_6$ & $v_7$ & $v_8$ & $v_9$ & $v_{10}$ \\
\midrule
$0$ (init) & $3$ & $4$ & $6$ & $6$ & $3$ & $4$ & $1$ & $1$ & $1$ & $1$ \\
$1$        & $3$ & $4$ & $6$ & $6$ & $3$ & $3$ & $0$ & $\boldsymbol{1}$ & $1$ & $1$ \\
$2$        & $3$ & $4$ & $6$ & $6$ & $3$ & $3$ & $\boldsymbol{1}$ &  & $1$ & $1$ \\
$3$        & $3$ & $3$ & $6$ & $6$ & $3$ & $3$ &  &  & $0$ & $\boldsymbol{1}$ \\
$4$        & $3$ & $3$ & $6$ & $6$ & $3$ & $3$ &  &  & $\boldsymbol{1}$ &  \\
$5$        & $1$ & $\boldsymbol{3}$ & $4$ & $4$ & $3$ & $3$ &  &  &  &  \\
$6$        & $\boldsymbol{3}$ &  & $3$ & $3$ & $3$ & $3$ &  &  &  &  \\
$7$        &  &  & $1$ & $\boldsymbol{3}$ & $1$ & $1$ &  &  &  &  \\
$8$        &  &  & $0$ &  & $\boldsymbol{3}$ & $0$ &  &  &  &  \\
$9$        &  &  & $\boldsymbol{3}$ &  &  & $0$ &  &  &  &  \\
$10$       &  &  &  &  &  & $\boldsymbol{3}$ &  &  &  &  \\
\bottomrule
\end{tabular}
\end{table}

\subsection{Static-\cpi{} implementation of change-only updates}\label{sec:changeonly-impl}

This subsection makes Algorithm~\ref{alg:peelH}'s abstract update $\supArr[u]\!\gets\!\supArr[u]-\Delta_v(u)$ concrete on the static \cpi{}.
%
The cliques that disappear from $u$'s support when $v$ leaves $R$ are exactly those encoded by a path $\TPath$ containing both $v$ and $u$ and alive in $G[R]$; the static \cpi{} therefore identifies the affected $u$ as the co-path neighbours of $v$, and evaluates each $\Delta_v(u)$ as a sum of per-path binomial differences.

\begin{lemma}[\cpi{} realisation of the change-only delta]\label{lem:changeonly-cpi}
Let $v\in R$ be popped and $u\in R\setminus\{v\}$.
%
Then
\[
  \Delta_v(u) \;=\; \sum_{\substack{\TPath\ni v,u \\ \TPath \text{ alive in } G[R]}} \pathDelta(\TPath, u, v),
\]
where for a path $\TPath$ alive in $G[R]$ with pre-pop pivot scope $p_{\mathrm{old}}=|\Vp(\TPath)\cap R|\ge\eta_{\TPath}$ and post-pop scope $p_{\mathrm{new}}=p_{\mathrm{old}}-\mathbf{1}[v\in\Vp(\TPath)]$, the per-path contribution is the binomial differential
\[
  \pathDelta(\TPath,u,v) \;=\;
  \begin{cases}
    \pathContrib(\TPath,u,q_{\mathrm{old}}), & v\in\Vh(\TPath) \text{ or } p_{\mathrm{new}}<\eta_{\TPath}, \\
    \pathContrib(\TPath,u,q_{\mathrm{old}}) - \pathContrib(\TPath,u,q_{\mathrm{new}}), & v\in\Vp(\TPath) \text{ and } p_{\mathrm{new}}\ge\eta_{\TPath},
  \end{cases}
\]
with $q_{\mathrm{old}}/q_{\mathrm{new}}$ following Definition~\ref{def:path-contrib}.
%
Each $\pathDelta$ evaluates in $O(1)$ via a precomputed binomial table, so applying all $\Delta_v(\cdot)$ updates for one pop costs $O\bigl(\sum_{\TPath\ni v}(|\Vh(\TPath)|+|\Vp(\TPath)|)\bigr)$.
\end{lemma}

\begin{proof}
By Lemma~\ref{lem:set-support-eval} applied at $R$ and $R\setminus\{v\}$,
\[
  \Delta_v(u)=\setSup_R(u)-\setSup_{R\setminus\{v\}}(u)=\sum_{\TPath\ni u}\bigl[W_{\TPath}^{R}(u)-W_{\TPath}^{R\setminus\{v\}}(u)\bigr],
\]
where $W_{\TPath}^{R}(u)=\mathbf{1}[\Vh(\TPath)\subseteq R]\cdot\pathContrib(\TPath,u,|\Vp(\TPath)\cap R|)$.
%
A path $\TPath$ with $v\notin\TPath$ has identical hold-inclusion and pivot scope in $R$ and $R\setminus\{v\}$, contributing $0$; the sum reduces to paths through both $v$ and $u$.
%
For such a path, $v\in\Vh(\TPath)$ makes $\mathbf{1}[\Vh(\TPath)\subseteq R\setminus\{v\}]=0$, killing the second term and giving the killed case.
%
For $v\in\Vp(\TPath)$, the hold inclusion is unaffected, and the pivot scope drops from $p_{\mathrm{old}}$ to $p_{\mathrm{new}}=p_{\mathrm{old}}-1$; the differential is $\pathContrib(\TPath,u,q_{\mathrm{old}})-\pathContrib(\TPath,u,q_{\mathrm{new}})$, with the second term vanishing by the binomial convention $\binom{m}{j}=0$ for $j>m$ when $p_{\mathrm{new}}<\eta_{\TPath}$.
%
The total cost claim follows because each $(\TPath,u)$ incidence with $\TPath\ni v$ is visited once per pop, and each visit triggers one $O(1)$ binomial lookup.
%
\qed
\end{proof}

In Algorithm~\ref{alg:peelH}, the iteration over ``$u\in R$ with $\Delta_v(u)\ne 0$'' (Line~7) is realised by walking every alive path $\TPath\ni v$ and visiting its co-path vertices; the per-path bookkeeping is the integer pivot counter $p_{\TPath}$ alongside the fixed $\eta_{\TPath}$, with the path alive iff $p_{\TPath}\ge\eta_{\TPath}$ (a derived predicate, with $p_{\TPath}\gets-1$ marking a hold-peeled path).
